% Part: first-order-logic
% Chapter: syntax-and-semantics
% Section: satisfaction

\documentclass[../../../include/open-logic-section]{subfiles}

\begin{document}

\olfileid{fol}{syn}{sat}

\olsection{Wanneer \article{formula} \printtoken{S}{formula}
  waar is in \article{structure} \printtoken{S}{structure}}

\begin{explain}
Termen en !!{formula}s horen bij !!{structure}s. Bij een term kijken we
naar zijn \emph{!!{value}}. Bij een formule kijken we of zij wordt
\emph{waar gemaakt}. De !!{value} van een term is een !!{element} van
!!a{structure}. Is de term alleen een constante, dan is zijn !!{value}
het object dat de !!{structure} aan die constante geeft. Is de term met
!!{function}s opgebouwd, dan rekenen we de !!{value} uit met de waarden
van de constanten en met de functies die in de structuur bij de
functietekens horen. !!^a{formula} wordt in !!a{structure} \emph{waar
gemaakt} als de uitleg van de predicaten de !!{formula} waar maakt voor
de objecten in het domein. Dit noemen we vervulling. We leggen stap
voor stap vast wanneer dat gebeurt. Bij een formule zonder logische
deelopbouw staat het antwoord al in de !!{structure}. Bij een langere
!!{formula} kijken we naar het buitenste logische teken, de
hoofdoperator, en naar de onmiddellijke !!{subformula}s.

Bij kwantoren is één extra stap nodig. De onmiddellijke !!{subformula}
van een gekwantificeerde !!{formula} bevat een vrije !!{variable}.
Een !!{structure} zegt niet welk object bij zo'n !!{variable} hoort.
Daarom gebruiken we ook \emph{toekenningen aan variabelen}. We zeggen
dan niet alleen in welke !!a{structure} een formule waar is, maar ook
bij welke !!a{variable} toekenning.
\end{explain}

\begin{defn}[Toekenning aan variabelen]
Een \emph{toekenning aan variabelen}, ook wel een bedeling genoemd,
is een functie~$s$ voor !!a{structure}~$\Struct{M}$. Ze koppelt elke
!!{variable} aan !!a{element} van~$\Domain M$. Dus:
$s\colon \Var \to \Domain M$.
\end{defn}

\begin{explain}
!!^a{structure} geeft elke !!{constant} een !!{value}. Een toekenning
aan variabelen geeft elke variabele een waarde. Ook een term die uit
zulke tekens is opgebouwd, moet een !!{element} van het !!{domain}
aanwijzen. Daarom bepalen we de !!{value} van een term stap voor stap.
Bij !!{constant}s en variabelen nemen we de waarde die de !!{structure}
of de toekenning al geeft. Bij een langere term rekenen we verder met
de functies die de !!{structure} aan de !!{function}s heeft gekoppeld.
\end{explain}

\begin{defn}[De !!^{value} van termen]
Neem een term $t$ van de taal~$\Lang L$, een !!{structure}~$\Struct M$
voor~$\Lang L$ en !!a{variable} toekenning~$s$ voor~$\Struct M$.
De \emph{!!{value}}~$\Value{t}{M}[s]$ bepalen we zo:
\begin{enumerate}
\item \indcase{t}{c}{$\Value{\indfrm}{M}[s] = \Assign{\indcomplex}{M}$.}
\item \indcase{t}{x}{$\Value{\indfrm}{M}[s] = s(\indcomplex)$.}
\item \indcase{t}{\Atom{f}{t_1, \ldots, t_n}}{
\[
\Value{\indfrm}{M}[s] = \Assign{f}{M}(\Value{t_1}{M}[s], \ldots,
\Value{t_n}{M}[s]).
\]}
\end{enumerate}
\end{defn}

\begin{defn}[$x$-variant]
Neem !!a{variable} toekenning~$s$ voor !!a{structure}~$\Struct M$.
Elke andere !!{variable} toekenning~$s'$ voor~$\Struct M$ die van~$s$
alleen mag verschillen in wat zij aan~$x$ geeft, heet een
\emph{$x$-variant} van~$s$. Als $s'$ zo'n $x$-variant van~$s$ is,
schrijven we $\varAssign{s'}{s}{x}$.
\end{defn}

\begin{explain}
Een $x$-variant van een toekenning~$s$ hoeft bij~$x$ \emph{niet}
werkelijk iets anders te geven. Elke toekenning is dus ook een
$x$-variant van zichzelf.
\end{explain}

\begin{defn}
  Neem !!a{variable} toekenning~$s$ voor !!a{structure}~$\Struct M$
  en neem $m \in \Domain{M}$. Dan is~$\Subst{s}{m}{x}$ de toekenning
  aan variabelen met de volgende regel:
  \[\Subst{s}{m}{x}(y) = \begin{cases}
    m & \text{als } y \ident x\\
    s(y) & \text{anders}.
  \end{cases}\]
\end{defn}

Kort gezegd is $\Subst{s}{m}{x}$ de ene $x$-variant van~$s$ die het
!!{element}~$m$ uit het domein aan~$x$ geeft. Alle andere
!!{variable}s dan~$x$ krijgen hetzelfde object als bij~$s$.

\begin{defn}[Wanneer een formule waar is]
\ollabel{defn:satisfaction}
We zeggen dat !!a{formula}~$!A$ waar is in !!a{structure}~$\Struct M$
bij !!a{variable} toekenning~$s$. We schrijven dat als
$\Sat{M}{!A}[s]$. De volgende regels leggen stap voor stap vast wanneer
dit geldt. (Met $\Sat/{M}{!A}[s]$ bedoelen we ``niet $\Sat{M}{!A}[s]$.'')
\begin{enumerate}
\tagitem{prvFalse}{%
  \indcase{!A}{\lfalse}{$\Sat/{M}{\indfrm}[s]$.}}{}

\tagitem{prvTrue}{%
  \indcase{!A}{\ltrue}{$\Sat{M}{\indfrm}[s]$.}}{}

\item \indcase{!A}{\Atom{R}{t_1, \dots, t_n}}{$\Sat{M}{\indfrm}[s]$
  precies als $\langle \Value{t_1}{M}[s], \dots, \Value{t_n}{M}[s] \rangle \in
  \Assign{R}{M}$.}

\item \indcase{!A}{\eq[t_1][t_2]}{$\Sat{M}{\indfrm}[s]$ precies als
  $\Value{t_1}{M}[s] = \Value{t_2}{M}[s]$.}

\tagitem{prvNot}{%
  \indcase{!A}{\lnot !B}{$\Sat{M}{\indfrm}[s]$ precies als
    $\Sat/{M}{!B}[s]$.}}{}

\tagitem{prvAnd}{%
  \indcase{!A}{(!B \land !C)}{$\Sat{M}{\indfrm}[s]$ precies als $\Sat{M}{!B}[s]$
    en $\Sat{M}{!C}[s]$.}}{}

\tagitem{prvOr}{%
  \indcase{!A}{(!B \lor !C)}{$\Sat{M}{\indfrm}[s]$ precies als
    $\Sat{M}{!B}[s]$ of $\Sat{M}{!C}[s]$ (of allebei).}}{}

\tagitem{prvIf}{%
  \indcase{!A}{(!B \lif !C)}{$\Sat{M}{\indfrm}[s]$ precies als $\Sat/{M}{!B}[s]$
    of $\Sat{M}{!C}[s]$ (of allebei).}}{}

\tagitem{prvIff}{%
  \indcase{!A}{(!B \liff !C)}{$\Sat{M}{\indfrm}[s]$ precies als ofwel zowel
    $\Sat{M}{!B}[s]$ als $\Sat{M}{!C}[s]$, ofwel noch $\Sat{M}{!B}[s]$
    noch $\Sat{M}{!C}[s]$.}}{}

\tagitem{prvAll}{%
  \indcase{!A}{\lforall[x][!B]}{$\Sat{M}{\indfrm}[s]$ precies als voor elk
    !!{element}~$m \in \Domain M$, $\Sat{M}{!B}[\Subst{s}{m}{x}]$.}}{}

\tagitem{prvEx}{%
  \indcase{!A}{\lexists[x][!B]}{$\Sat{M}{\indfrm}[s]$ precies als voor minstens
  één !!{element}~$m \in \Domain M$, $\Sat{M}{!B}[\Subst{s}{m}{x}]$.}}{}
\end{enumerate}
\end{defn}

\begin{explain}
In de laatste
\iftag{notprvEx,notprvAll}{regel}{twee regels} hebben we de toekenningen
aan variabelen echt nodig.\iftag{prvAll}{ We kunnen niet zeggen dat
$\lforall[x][!B(x)]$ waar is als ``voor alle $m \in \Domain{M}$,
$\Sat{M}{!B(m)}$.''}{}\iftag{prvEx}{ We kunnen niet zeggen dat
$\lexists[x][!B(x)]$ waar is als ``voor minstens één $m \in \Domain{M}$,
$\Sat{M}{!B(m)}$.''}{} Een $m \in \Domain M$ is namelijk geen teken
van de taal. Daarom is $!B(m)$ geen !!a{formula}; anders gezegd is
$\Subst{!B}{m}{x}$ niet gedefinieerd. Ook mogen we er niet van uitgaan
dat er genoeg !!{constant}s of termen zijn om elk !!{element} van
~$\Struct{M}$ een naam te geven. De definitie van !!{structure}s eist
dat niet. In de standaardtaal zijn er !!{denumerable} veel
!!{constant}s. Als $\Domain{M}$ niet !!{enumerable} is, zijn er dus te
weinig !!{constant}s om ieder object te benoemen.

De oplossing is een !!{variable} toekenning. Daarmee koppelen we een
variabele rechtstreeks aan een !!{element} van het domein. We zeggen
dus niet bijvoorbeeld dat
\iftag{prvEx}{$\lexists[x][!B(x)]$}{$\lforall[x][!B(x)]$} waar is in
~$\Struct M$ precies als \iftag{prvEx}{voor minstens één $m \in
\Domain{M}$ de voorwaarde bij de formule geldt}{voor alle $m \in \Domain{M}$, $\Sat{M}{!B(m)}$}. We zeggen dat zij waar is in~$\Struct M$
\emph{bij}~$s$ precies als $!B(x)$ waar is bij~$\Subst{s}{m}{x}$,
\iftag{prvEx}{voor minstens één}{voor elk} $m \in \Domain M$.
\end{explain}

\begin{ex}
Neem $\Lang{L} = \{a, b, f, R\}$. Hier zijn $a$ en $b$ !!{constant}s,
$f$~is een functie met twee invoeren en $R$~een tweeplaatsig
!!{predicate}. Neem de volgende !!{structure}~$\Struct{M}$:
\begin{enumerate}
\item $\Domain M = \{1, 2, 3, 4\}$
\item $\Assign{a}{M} = 1$
\item $\Assign{b}{M} = 2$
\item $\Assign{f}{M}(x, y) = x+y$ als $x+y \le 3$ en anders $= 3$.
\item $\Assign{R}{M} = \{\tuple{1, 1}, \tuple{1, 2}, \tuple{2, 3}, \tuple{2, 4}\}$
\end{enumerate}
De functie $s(x) = 1$ geeft elke !!{variable} het element
$1 \in \Domain{M}$. Zij is dus een toekenning aan variabelen voor
~$\Struct{M}$.

Nu rekenen we:
\begin{align*}
\Value{f(a,b)}{M}[s] & = \Assign{f}{M}(\Value{a}{M}[s], \Value{b}{M}[s]).
\intertext{$a$ en $b$ zijn !!{constant}s. Daarom is $\Value{a}{M}[s]
  = \Assign{a}{M} = 1$ en $\Value{b}{M}[s] = \Assign{b}{M} = 2$. Dus}
\Value{f(a,b)}{M}[s] & = \Assign{f}{M}(1, 2) = 1+2 = 3.
\intertext{Voor de waarde van $f(f(a,b),a)$ bekijken we eerst}
\Value{f(f(a,b),a)}{M}[s] & = \Assign{f}{M}(\Value{f(a, b)}{M}[s],
\Value{a}{M}[s]) = \Assign{f}{M}(3, 1) = 3,
\intertext{want $3+1 > 3$. Verder is $s(x) = 1$ en $\Value{x}{M}[s] =
  s(x)$. Daarom is ook}
\Value{f(f(a,b),x)}{M}[s] & = \Assign{f}{M}(\Value{f(a, b)}{M}[s],
\Value{x}{M}[s]) = \Assign{f}{M}(3, 1) = 3,
\end{align*}

Een formule zonder logische deelopbouw, zoals $R(t_1, t_2)$, is waar
als het paar waarden van haar argumenten in de relatie staat. Dat paar
is $\tuple{\Value{t_1}{M}[s], \Value{t_2}{M}[s]}$ en moet !!a{element}
zijn van~$\Assign{R}{M}$. Daarom geldt $\Sat{M}{R(b,f(a,b))}[s]$: we
hebben $\tuple{\Value{b}{M}, \Value{f(a,b)}{M}} = \tuple{2, 3} \in
\Assign{R}{M}$. Maar $\Sat/{M}{R(x, f(a,b))}[s]$, want
$\tuple{1, 3} \notin \Assign{R}{M}$.

Bij een langere formule~$!A$ kijk je eerst naar het buitenste logische
teken, de hoofdoperator. Daarna gebruik je de regel die daarbij hoort.
In $R(a, a) \lif (R(b, x) \lor R(x, b))$ is de hoofdoperator~$\lif$. Dus
\begin{align*}
  & \Sat{M}{R(a, a) \lif (R(b, x) \lor R(x, b))}[s] \text{ precies als }\\
  & \qquad
  \Sat/{M}{R(a,a)}[s] \text{ of } \Sat{M}{R(b, x) \lor R(x, b)}[s]
  \intertext{$\Sat{M}{R(a,a)}[s]$ geldt, want $\tuple{1,1} \in
    \Assign{R}{M}$. We weten het antwoord dus nog niet. Eerst moeten we
    kijken of $\Sat{M}{R(b, x) \lor R(x, b)}[s]$:}
  & \Sat{M}{R(b, x) \lor R(x, b)}[s] \text{ precies als }\\
  & \qquad \Sat{M}{R(b,x)}[s] \text{ of } \Sat{M}{R(x, b)}[s]
  \intertext{Dat klopt, want $\Sat{M}{R(x, b)}[s]$ en
    $\tuple{1,2} \in \Assign{R}{M}$.}
\end{align*}

Een $x$-variant van~$s$ is een toekenning die van $s$ alleen mag
verschillen in wat zij aan~$x$ geeft. Voor elk !!{element} van
~$\Domain{M}$ hebben we zo'n $x$-variant van~$s$:
\begin{align*}
  s_1 & = \Subst{s}{1}{x}, &
  s_2 & = \Subst{s}{2}{x},\\
  s_3 & = \Subst{s}{3}{x}, &
  s_4 & = \Subst{s}{4}{x}.
\end{align*}
Zo is $s_2(x) = 2$ en $s_2(y) = s(y) = 1$ voor elke variabele~$y$
behalve~$x$. Dit zijn alle $x$-varianten van~$s$ voor
~$\Struct{M}$, want $\Domain{M} = \{1, 2, 3, 4\}$. Let vooral op dat
$s_1 = s$: $s$~is altijd een $x$-variant van zichzelf.

\iftag{prvEx}{Om te weten of de existentieel gekwantificeerde
  !!{formula}~$\lexists[x][!A(x)]$ waar is, moeten we nagaan of
  $\Sat{M}{!A(x)}[\Subst{s}{m}{x}]$ geldt voor minstens één $m \in \Domain M$. Dus
  \[
  \Sat{M}{\lexists[x][(R(b,x) \lor R(x,b))]}[s],
  \]
  want $\Sat{M}{R(b,x) \lor R(x, b)}[\Subst{s}{1}{x}]$ geldt.
  ($\Subst{s}{3}{x}$ zou ook werken.) Maar
  \[
  \Sat/{M}{\lexists[x][(R(b,x) \land R(x,b))]}[s]
  \]
  want voor elke keuze van $m \in \Domain{M}$ geldt $\Sat/{M}{R(b,x) \land R(x,b)}[\Subst{s}{m}{x}]$.}{}

\iftag{prvAll}{Om te weten of de universeel gekwantificeerde
  !!{formula}~$\lforall[x][!A(x)]$ waar is, moeten we nagaan of
  $\Sat{M}{!A(x)}[\Subst{s}{m}{x}]$ geldt voor elke $m \in \Domain M$. Dus
  \[
  \Sat{M}{\lforall[x][(R(x,a) \lif R(a,x))]}[s],
  \]
  want $\Sat{M}{R(x,a) \lif R(a,x)}[\Subst{s}{m}{x}]$ geldt voor elke $m \in
  \Domain M$. Bij $m = 1$ geldt $\Sat{M}{R(a,x)}[\Subst{s}{1}{x}]$,
  dus is het tweede deel waar. Bij $m = 2$, $3$ en~$4$ geldt
  $\Sat/{M}{R(x,a)}[\Subst{s}{m}{x}]$, dus is het eerste deel onwaar.
  Maar
  \[
  \Sat/{M}{\lforall[x][(R(a,x) \lif R(x,a))]}[s]
  \]
  want $\Sat/{M}{R(a,x) \lif R(x,a)}[\Subst{s}{2}{x}]$: hier geldt
  $\Sat{M}{R(a, x)}[\Subst{s}{2}{x}]$ en $\Sat/{M}{R(x,
  a)}[\Subst{s}{2}{x}]$.}{}

\iftag{defEx}{Om te weten of de existentieel gekwantificeerde
  !!{formula}~$\lexists[x][!A(x)]$ waar is, kijken we of
  $\Sat{M}{\lnot \lforall[x][\lnot !A(x)]}[s]$. Bijvoorbeeld:
  \[
  \Sat{M}{\lexists[x][(R(b,x) \lor R(x,b))]}[s].
  \]
  Eerst geldt $\Sat{M}{R(b,x) \lor R(x, b)}[\Subst{s}{1}{x}]$.
  ($\Subst{s}{3}{x}$ zou ook werken.) Daarom geldt
  $\Sat/{M}{\lnot(R(b,x) \lor R(x,b))}[\Subst{s}{1}{x}]$, en dus
  $\Sat/{M}{\lforall[x][\lnot(R(b,x) \lor R(x,b))]}[s]$. Daarmee geldt
  $\Sat{M}{\lnot\lforall[x][\lnot(R(b,x) \lor R(x,b))]}[s]$. Maar
  \[
  \Sat/{M}{\lexists[x][(R(b,x) \land R(x,b))]}[s].
  \]
  Dit komt doordat $\Sat{M}{\lforall[x][\lnot(R(b,x) \land
  R(x,b))]}[s]$. Er is namelijk geen $m \in \Domain M$ waarvoor
  $\Sat{M}{R(b,x) \land R(x,b)}[\Subst{s}{m}{x}]$. Deze voorbeelden
  laten de algemene regel al zien: $\Sat{M}{\lexists[x][!A(x)]}[s]$
  precies als $\Sat{M}{!A(x)}[\Subst{s}{m}{x}]$ geldt voor minstens één
  $m \in \Domain M$.}{}

\iftag{defAll}{Om te weten of de universeel gekwantificeerde
  !!{formula}~$\lforall[x][!A(x)]$ waar is, kijken we of
  $\Sat{M}{\lnot\lexists[x][\lnot !A(x)]}[s]$. Bijvoorbeeld:
  \[
  \Sat{M}{\lforall[x][(R(x,a) \lif R(a,x))]}[s],
  \]
  Eerst geldt $\Sat{M}{R(x,a) \lif R(a,x)}[\Subst{s}{m}{x}]$ voor elke $m \in
  \Domain M$. Hier geldt $\Sat{M}{R(a,x)}[\Subst{s}{1}{x}]$, en
  $\Sat/{M}{R(x,a)}[\Subst{s}{m}{x}]$ bij $m = 2$, $3$ of~$4$.
  Er is dus geen $m \in \Domain M$ waarvoor
  $\Sat{M}{\lnot(R(x,a) \lif R(a,x))}[\Subst{s}{m}{x}]$. Daarom geldt
  $\Sat/{M}{\lexists[x][\lnot(R(x,a) \lif R(a,x))]}[s]$, en dus
  $\Sat{M}{\lnot\lexists[x][\lnot(R(x,a) \lif R(a,x))]}[s]$. Maar
  \[
  \Sat/{M}{\lforall[x][(R(a,x) \lif R(x,a))]}[s],
  \]
  want $\Sat/{M}{R(a,x) \lif R(x,a)}[\Subst{s}{2}{x}]$. Daarom geldt
  $\Sat{M}{\lexists[x][\lnot(R(a,x) \lif R(x,a))]}[s]$. De algemene
  regel is dus: $\Sat{M}{\lforall[x][!A(x)]}[s]$ precies als
  $\Sat{M}{!A(x)}[\Subst{s}{m}{x}]$ geldt voor elke $m \in \Domain M$.}{}

Nu een ingewikkelder voorbeeld:
\[
\lforall[x][(R(a,x) \lif \lexists[y][R(x,y)])].
\]
We hebben $\Sat/{M}{R(a,x)}[\Subst{s}{3}{x}]$ en
$\Sat/{M}{R(a,x)}[\Subst{s}{4}{x}]$. Daarom hoeven we het tweede deel
van de implicatie alleen te bekijken bij $m = 1$ en $m = 2$. Geldt
$\Sat{M}{\lexists[y][R(x,y)]}[\Subst{s}{1}{x}]$? Dat is zo als er
minstens één $n \in \Domain M$ is waarvoor
$\Sat{M}{R(x,y)}[\Subst{\Subst{s}{1}{x}}{n}{y}]$. Neem $n = 1$.
Dan is $\Subst{\Subst{s}{1}{x}}{n}{y} = \Subst{s}{1}{y} = s$.
Verder is $s(x) = 1$, $s(y) = 1$ en $\tuple{1,1} \in \Assign{R}{M}$.
Het antwoord is dus ja.

Voor $\Sat{M}{\lexists[y][R(x,y)]}[\Subst{s}{2}{x}]$ bekijken we de
!!{variable} toekenningen $\Subst{\Subst{s}{2}{x}}{n}{y}$. Bij
$n = 1$ krijgen we~$s_2 = \Subst{s}{2}{x}$. Die maakt $R(x,y)$ niet
waar: $s_2(x) = 2$, $s_2(y) = 1$ en $\tuple{2,1}\notin \Assign{R}{M}$.
Neem nu $\Subst{\Subst{s}{2}{x}}{3}{y} = \Subst{s_2}{3}{y}$.
Dan geldt $\Sat{M}{R(x,y)}[\Subst{s_2}{3}{y}]$, want
$\tuple{2,3} \in \Assign{R}{M}$. Dus geldt
$\Sat{M}{\lexists[y][R(x,y)]}[s_2]$.

Voor elke $m \in \Domain M$ geldt nu één van twee dingen:
$\Sat/{M}{R(a,x)}[\Subst{s}{m}{x}]$ (als $m = 3$, $4$), of
$\Sat{M}{\lexists[y][R(x,y)]}[\Subst{s}{m}{x}]$ (als $m = 1$, $2$). Daarom
\[
\Sat{M}{\lforall[x][(R(a,x) \lif \lexists[y][R(x,y)])]}[s].
\]
Maar
\[
\Sat/{M}{\lexists[x][(R(a,x) \land \lforall[y][R(x,y)])]}[s].
\]
$\Sat{M}{R(a,x)}[\Subst{s}{m}{x}]$ geldt alleen voor $m = 1$ en $m =
2$. Bij allebei die waarden van~$m$ bestaat er een $n \in \Domain M$
waarvoor die formule niet waar is. Bij de eerste waarde nemen we $n = 4$;
bij de tweede nemen we het domeinelement 1. Dan geldt
$\Sat/{M}{R(x,y)}[\Subst{\Subst{s}{m}{x}}{n}{y}]$, en dus ook
$\Sat/{M}{\lforall[y][R(x,y)]}[\Subst{s}{m}{x}]$ voor $m = 1$ en $m =
2$. Er is dus geen $m \in \Domain M$ waarvoor
$\Sat{M}{R(a,x) \land \lforall[y][R(x,y)]}[\Subst{s}{m}{x}]$.
\end{ex}

\iftag{defEx}{%
\begin{prop}\ollabel{prop:sat-ex}
$\Sat{M}{\lexists[x][!B(x)]}[s]$ precies als er een $x$-variant $s'$ van $s$
    bestaat waarvoor $\Sat{M}{!B(x)}[s']$.
\end{prop}

\begin{proof}
  Oefening.
\end{proof}
}{}

\tagprob{defEx}
\begin{prob}
  Bewijs \olref[fol][syn][sat]{prop:sat-ex}
\end{prob}
\tagendprob

\iftag{defAll}{%
\begin{prop}\ollabel{prop:sat-all}
$\Sat{M}{\lforall[x][!B(x)]}[s]$ precies als voor elke $x$-variant~$s'$ van $s$
  geldt dat $\Sat{M}{!B(x)}[s']$
\end{prop}

\begin{proof}
  Oefening.
\end{proof}
}{}

\tagprob{defAll}
\begin{prob}
  Bewijs \olref[fol][syn][sat]{prop:sat-all}
\end{prob}
\tagendprob

\begin{prob}
Neem $\Lang L = \{c, f, A\}$ met één !!{constant}, één !!{function}
met één invoer en één !!{predicate} met twee plaatsen. Neem de volgende
!!{structure}~$\Struct{M}$:
\begin{enumerate}
\item $\Domain M = \{1, 2, 3\}$
\item $\Assign{c}{M} = 3$
\item $\Assign{f}{M}(1) = 2, \Assign{f}{M}(2) = 3, \Assign{f}{M}(3) = 2$
\item $\Assign{A}{M} = \{\tuple{1, 2}, \tuple{2, 3}, \tuple{3, 3}\}$
\end{enumerate}
(a) Laat $s(v) = 1$ voor elke !!{variable}~$v$. Zoek uit of
\[
\Sat{M}{\lexists[x][(A(f(z), c) \lif \lforall[y][(A(y, x) \lor A(f(y),
      x))])]}[s]
\]
geldt. Leg uit waarom wel of niet.

(b) Geef een andere structuur en !!a{variable} toekenning waarin de
!!{formula} niet waar is.
\end{prob}

\end{document}
